\section{Implementation}
\label{sec:implementation}

\subsection{Architecture and Reproducibility}
\label{sec:architecture}

The prototype is written in Python 3.11. Dependencies are pinned in
\texttt{requirements.txt}, and the random seed and configuration file are committed with the
code. The core routing logic is a single small module; the review interface is a command-line
tool, so the workflow can be reproduced on a server without a graphical environment. Each
stage in Figure~\ref{fig:pipeline} has its own test, and the audit trail is written as JSON
Lines, one action per line.

The function below is a minimal example of the routing logic: a document passes automatically
when the model score is above the threshold and an evidence span exists; otherwise it goes to
human review.

\inputminted[fontsize=\footnotesize]{python}{code/example.py}

\subsection{Components and Responsibilities}
\label{sec:components}

Table~\ref{tab:components} lists the main components and their responsibilities. The
interfaces are explicit: the classifier returns only a label and a score, the claim extractor
returns only text spans and offsets, and the audit trail is append-only.

\begin{table}[htbp]
  \centering
  \formaltable
  \begin{tabularx}{\textwidth}{|B|Y|}
    \hline
    \textbf{Component} & \textbf{Responsibility} \\ \hline
    Document intake & Reads the document and metadata; checks format and authorisation fields. \\ \hline
    Classifier & Proposes a topic label and a confidence score; never rejects a document. \\ \hline
    Claim extractor & Returns the text spans and offsets that support the label. \\ \hline
    Human review & Accepts, edits or rejects the proposal; the final label is human. \\ \hline
    Audit trail & Appends model output, review action, time and reason. \\ \hline
  \end{tabularx}
  \caption{System components and responsibilities.}
  \label{tab:components}
\end{table}

